\section{Related Work}
\paragraph{Image manipulation localization.}
In a locally edited photograph, most pixels may remain authentic. TruFor combines RGB features with a learned noise representation and uses deviations from expected image statistics to localize a manipulation~\cite{trufor}. DeCLIP studies the complementary question of how much spatial evidence a decoder can recover from pretrained CLIP features~\cite{declip}. These approaches motivate using both scene content and local residuals. A residual response may also arise from genuine texture or image processing, so its contribution must be assessed on genuine images as well as edited ones.

\paragraph{Diffusion inpainting.}
OpenSDI introduced OpenSDID to evaluate diffusion image detection and localization across generators and manipulation types. Its MaskCLIP method combines CLIP and masked autoencoder representations~\cite{opensdi}. FUSED combines a frozen semantic encoder and a forensic encoder through token routing and reports results for several methods under an SD1.5-only training protocol~\cite{fused}. We use that evaluation to contextualize prior work. Our local experiments examine how whole-image and crop-level supervision affect false alarms and missed edits on a development subset with a different training and evaluation protocol.

CAT-Net uses JPEG compression evidence, PSCC-Net predicts masks progressively across scales, and ObjectFormer combines RGB and frequency features through an object-based transformer~\cite{catnet,psccnet,objectformer}. AutoSplice provides a complementary benchmark for text-guided local edits~\cite{autosplice}. We leave evaluation on that benchmark to future work.

\paragraph{Evidence outside the edited region.}
Latent diffusion operates in a learned image representation, and editors often share components with systems used for full-image synthesis~\cite{ldm,sdxl}. Earlier detectors exploit convolutional artifacts, texture, frequency information, or neighboring-pixel structure~\cite{cnnspot,gramnet,freqnet,npr}. Other approaches learn classification from pretrained visual features~\cite{clip,unifd,rine} or use reconstruction errors~\cite{dire,aeroblade}. These cues can support image classification without specifying where the edit occurred. The inpainting exchange study demonstrates the risk of relying on an image-wide trace when the underlying modification is local~\cite{inpx}. This motivates recording a localization outcome alongside image-level accuracy.

\paragraph{Spatial evidence and evaluation.}
Manipulation tracing, contrastive localization, uncertainty-guided learning, and multi-view models provide alternatives for connecting a region prediction to image evidence~\cite{mantranet,cfl,uncertainty,mvss}. IML-ViT and SparseViT study transformer representations for manipulation localization, while INPDiffusion targets inpainting masks directly~\cite{imlvit,sparsevit,inpdiff}. In our experiments, the question is whether local predictions help recover missed edits at a usable false-positive rate. We therefore retain counts for both classes and report localization by edit size.
